\cotefile{77}
\scanpage{114}
Soit donc donné (\unsure{après une|après|supposons} extension peut-être de $S$ en $S[\underline{i}]\ldots$) $w=iw_0$ de sorte que $\delta'(q)=w^{\otimes2}$, il faut calculer l'effet de $g_q$ sur $u+\underline{i}v$, i.e. de $f_w$ sur $u+iv,u-iv$.

La \unsure{recette|vérité} est la suivante : On considère $(E,-q,w=i\unsure{)}\unsure{w|\omega}_0)\simeq(\unsure{\gamma|\mathfrak g}M',q_{M'},w_{M'})$ et on transforme $u+iv,u-iv$ en \unsure{élém\textsuperscript{ts}|éléments} de $\gamma M\ldots$

\unsure{OBS|OPS} \quad $(E,q,w_0)=(\gamma M,q_M,w_M)$,
\inserted{i.e. $(E,-q,w)=(\gamma M,-q_M,iw_M)$} et on peut prendre $M'=M$, $\gamma M'\longrightarrow\gamma M$ donné par $\theta=i\,\mathrm{id}_{\gamma M}$, qui transforme $(u+iv,u-iv)$ en
\[
\begin{aligned}
\alpha&=iu-v\\
\beta&=iu\deleted{-}+v
\end{aligned}
\]
d'où
\[
\begin{aligned}
2iu&=\alpha+\beta\\
2v&=-\alpha+\beta
\end{aligned}
\]
\[
\left\{\begin{aligned}
4i[u,v]&=2[\alpha,\beta]\qquad\text{i.e.}\quad 2i[u,v]=[\alpha,\beta]\\
-4q(u)&=\deleted{q(\alpha)}+\deleted{q(\beta)}+\varphi(\alpha,\beta)=\operatorname{tr}\alpha\beta\\
4q(\unsure{v|u})&=q(\alpha)+q(\beta)-\varphi(\alpha,\beta)=-\operatorname{tr}\alpha\beta
\end{aligned}\right.
\]
\[
[u,v]/2q(u)=\frac{1}{2i}\frac{-\unsure{2}}{\operatorname{tr}\alpha\beta}[\alpha\unsure{,}\beta]=\frac{i}{\operatorname{tr}\alpha\beta}[\alpha,\beta\unsure{)|]}\quad\text{i.e.}
\]
i.e.
\[
\theta\bigl([u,v]/2q(u)\bigr)=i\frac{[u,v]}{2q(u)}=-[\alpha,\beta]/\operatorname{tr}\alpha\beta
\]
% archive sheet number: 113
\scanpage{115}
On trouve l'identité \deleted{\unsure{des|dès qu'}} on trouve comme \unsure{isom|isom.} $g_M$
\[
X_M\times X_M\longrightarrow\widehat{\Sigma}_M
\]
celui qui, pour $\alpha,\beta\in\Gamma\unsure{\gamma|\mathfrak g}M$, définissant des sections de $X$ (i.e. tels que $\alpha,\beta$ aient tr et det nuls, et sont $\ne0$ sur tte fibre) distinctes en chaque fibre (i.e.
\[
\operatorname{tr}(\alpha\beta)\ne0\qquad\text{sur chaque fibre),}
\]
associe la \deleted{élément}\inserted{section} de $\Sigma\subset\inserted{\gamma M}$ définie par
\[
-[\alpha,\beta]/\unsure{\operatorname{tr}|\varphi}(\alpha\unsure{,}\beta)
\]
(et non $+[\alpha,\beta]/\operatorname{tr}\alpha\beta$). Il faut vérifier sur mes propres textes \unsure{canoniques|cosmiques|conséquente} \ldots. Si on veut l'interpréter comme section de $\widehat\Sigma$, il faut prendre $-[\alpha,\beta]/\operatorname{tr}(\alpha\beta)+1$ ou encore $\operatorname{tr}(\alpha\beta)\cdot1-[\alpha,\beta]$. Donc l'iso \deleted{$\unsure{X_M\times X}$}
\[
X_M\times X_M\unsure{-|\longrightarrow}\widehat{\Sigma}_M
\]
\deleted{semble} provient d'un \deleted{iso} \inserted{hom} \deleted{\unsure{sur}} au niveau des cônes \deleted{\unsure{épointés|associés}} affines $C=\widetilde V_{\mathrm{aff}}(q_{\gamma M})$ et $D=V_{\mathrm{aff}}(q_M)$ \deleted{$\widetilde X_M\times\widetilde X_M\longrightarrow\unsure{\widetilde\Sigma}$}
% archive sheet number: 114
\scanpage{116}
\[
\deleted{(\alpha,\beta)\longmapsto}
\]
\[
\begin{aligned}
F_M&:C_M\times C_M\unsure{-|\longrightarrow}D_M\\
F_M(\alpha,\beta)&=\operatorname{tr}(\alpha\beta).1-[\alpha,\beta]
\end{aligned}
\]
Compatible avec l'hom
\[
G_m\times G_m\xrightarrow{(\lambda,\mu)\unsure{-|\mapsto}\lambda\mu}G_m
\]
sur les groupes d'homothéties. \deleted{\unsure{Passant aux} cônes épointés, cela se traduit par un hom de $G_m$-torseurs}
\[
\deleted{C_M^{\unsure{\times|\wedge}}\times C_M^{\unsure{\times|\wedge}}/\underset{\text{\unsure{agissant|opérant|épointé} antidiagonalement}}{G_m}\ \underset{\unsure{-}}{\sim}\ D_M^{\unsure{\times|\wedge}}}
\]

Le fait que l'on s'envoie bien dans $D_M$ provient de la formule
\[
\begin{aligned}
&\alpha,\beta\ \text{\unsure{str.|étr.|étant} nilpotents}\qquad(\text{i.e. }\det\alpha=\det\beta=\operatorname{tr}\alpha=\operatorname{tr}\beta=0)\\
&\Longrightarrow q_M\bigl([\alpha,\beta]/\varphi(\alpha,\beta)\bigr)=1\\
&\hspace{2em}\text{i.e.}\quad\det([\alpha,\beta])+\unsure{\operatorname{tr}|\varphi}(\alpha\beta)^2=0
\end{aligned}
\]
i.e.
\[
q_M[\alpha,\beta]=\varphi(\alpha,\beta{\unsure{]|)}}^2
\]
ce qui provient en effet de la formule générale (valable pour $\alpha,\beta\in\Gamma\unsure{\gamma|\mathfrak g}M$ sans plus)
\[
q_M\bigl([\alpha,\beta]\bigr)=-4q(\alpha)q(\beta)+\varphi(\alpha,\beta)^2
\]
% archive sheet number: 115
\scanpage{117}
\[
\det([\alpha,\beta])+\bigl(\operatorname{tr}(\alpha\beta)\bigr)^2-4(\det\alpha)(\det\beta)=0\tag*{(0)}
\]
si $\operatorname{tr}\alpha=\operatorname{tr}\beta=0$

Soit $u=\alpha+s.1,\ v=\beta+t.1$ donc $\begin{aligned}\operatorname{tr}u&=2s\\\operatorname{tr}v&=2t\end{aligned}$

d'où
\[
[u,v]=[\alpha,\beta]\tag*{(\unsure{1|*|4})}
\]
\[
\begin{aligned}
&\deleted{\operatorname{tr}(\alpha\beta)=\operatorname{tr}(u-s)(v-t)=\operatorname{tr}(uv)+2st}\\
&\deleted{\operatorname{tr}\alpha\beta=\operatorname{tr}(u-s)(v-t)=\operatorname{tr}uv-\underbrace{s\operatorname{tr}v}_{1/2\,\operatorname{tr}u\operatorname{tr}v}-\underbrace{t\operatorname{tr}u}_{1/2\,\operatorname{tr}v\operatorname{tr}u}+\underbrace{2st}_{\unsure{1/2\,\operatorname{tr}u|1/2\,\operatorname{tr}u\operatorname{tr}v}}}\\
&\deleted{{}=\operatorname{tr}uv-\operatorname{tr}u\operatorname{tr}v}
\end{aligned}
\]
\[
\begin{aligned}
uv&=\alpha\beta+s\deleted{\unsure{\alpha}}\beta+t\deleted{\unsure{\beta}}\alpha+st\qquad\text{donc}\\
\operatorname{tr}uv&=\operatorname{tr}\alpha\beta+2st=\operatorname{tr}\alpha\beta+\tfrac12\operatorname{tr}u\operatorname{tr}v
\end{aligned}
\]
i.e.
\[
\operatorname{tr}\alpha\beta=\operatorname{tr}uv-\tfrac12\operatorname{tr}u\operatorname{tr}v
\]
\[
(\operatorname{tr}\alpha\beta)^2=(\operatorname{tr}uv)^2-(\operatorname{tr}uv)(\operatorname{tr}u)(\operatorname{tr}v)+\tfrac14(\operatorname{tr}u)^{2}(\operatorname{tr}v)^2\tag*{(2)}
\]
\[
\left\{\begin{aligned}
\det\alpha&=\det(u-s)=\det(s-u)=s^2+\det u=\tfrac14(\operatorname{tr}u)^2+\det\unsure{u}\\
\det\beta&=t^2+\det v=\tfrac14(\operatorname{tr}v)^2+\det v
\end{aligned}\right.\tag*{(3)}
\]
donc
\[
\begin{aligned}
\det([u,v])+(\operatorname{tr}uv)^2&-(\operatorname{tr}uv)(\operatorname{tr}u)(\operatorname{tr}v)+\tfrac14\overline{(\operatorname{tr}u)^2(\operatorname{tr}v)^2}\\
&-\tfrac14[\deleted{\unsure{s^2}}(\operatorname{tr}u)^2+4\det u][(\operatorname{tr}v)^{\unsure{2}}+4\det v]
\end{aligned}
\]
\[
\boxed{\begin{aligned}
\det([u,v])&\unsure{-|+}(\det u)(\operatorname{tr}v)^2-(\det v)(\operatorname{tr}u)^2-4(\det u)(\det v)\\
&+(\operatorname{tr}uv)^2-(\operatorname{tr}uv)(\operatorname{tr}u)(\operatorname{tr}v)=0
\end{aligned}}
\]
% archive sheet number: 116
\scanpage{118}
\deleted{\fbox{1''}} Quadriques affines \inserted{bi}régulières (i.e. \inserted{\unsure{project\textsuperscript{t}|projectivement}} lisses et lisses à l'infini \unsure{voir --|non --})

1. Soient, sur schéma $S$, $P$ un fibré projectif de rang $2\nu-1$, $Q$ une quadrique lisse dans $P$, $H$ un hyperplan de $P$ non tangent à $Q$, $\unsure{P|Q}-H=E$ le fibré affine correspondant. Donc $E$ (ou encore $(P,H)$) est déterminé par une extension de $\mathcal O_S$-fibrés vect.
\diagram{\[
\begin{tikzcd}[ampersand replacement=\&,column sep=large]
0 \arrow[r] \& V \arrow[r] \& T \arrow[r,"\pi"] \& \mathcal O_S \arrow[r,phantom,"{\unsure{-|\longrightarrow}}" description] \& 0
\end{tikzcd}\tag*{(1)}
\]}
où $\underline V$ est le $\mathcal O_S$-Module loc. libre (de rang $2\nu-1$) des translations de $T$, $E=\pi^{-1}(1)$, \inserted{$T$ étant de rg $2\nu$.} Quant à $Q$, elle est déterminée par la donnée d'une section de $\check{\mathbb P}(\operatorname{Sym}^2(\check T))=\mathbb P(\deleted{\unsure{\operatorname{Sym}}}\Gamma^2(T))$, i.e. une forme quadratique \inserted{$q$} sur $T$ définie loc. mod. multiplication par une section inversible.

La condition $Q$ lisse signifie $q$ lisse i.e. la forme bil. associée $\varphi$ non dégénérée, ou encore que
\[
\delta(q)\in\Gamma(S,(\det T)^{\otimes-2})\tag*{(2)}
\]
est une \deleted{section}\inserted{base} de $(\underline{\det}\,T)^{\otimes-2}$, i.e. est partout non \unsure{nulle}. La condition que $H$ est non tangent à $Q$ signifie aussi que $q\mid V$ est une forme lisse. \deleted{ou encore que la droite ortho}\inserted{On va la} reformuler, en notant $D$ le \unsure{$k$-fibré|fibré|ss-fibré|sous-fibré} inversible
% archive sheet number: 117
\scanpage{119}
de $T$ défini par
\[
\unsure{D|\mathbb D}=V^\perp\tag*{(3)}
\]
et en notant que
\[
D\simeq(E/V)^\vee\simeq\mathcal O_S^\vee\simeq\mathcal O_S\tag*{(4)}
\]
i.e. $D$ est muni canoniquement d'une \deleted{section}\inserted{base} $a$, \unsure{d.à.d.|i.e.} d'une section
\[
\left\{\begin{aligned}
&a\in\Gamma(S,T)\deleted{\unsure{\subset\Gamma(|\subset\Gamma(T)}}\\
&E=\{x\in\Gamma(-,T)\mid\varphi(x,a)=1\}
\end{aligned}\right.\tag*{(4')}
\]
-- du moins quand $q$ est \uline{\unsure{choisie|définie|donnée}} \deleted{\unsure{sans}} (non seulement mod une unité). Dans le cas général, regardant $q$ comme forme à valeurs dans un module inv. $L$ (quotient de $\Gamma^2(T)$, i.e. dual d'un ss-fibré inv. de $\operatorname{Sym}^2(\check T)\ldots$) on trouve $D\otimes(\unsure{T|E}/V)\simeq L$ i.e. [comme $T/V\simeq\mathcal O_S$] on trouve un isom canonique
\[
D\xrightarrow{\sim}L\tag*{(5)}
\]
La condition $H$ non tangent à $Q$ s'exprime alors par la condition que
% archive sheet number: 118
\scanpage{120}
$q\mid D$ non nul fibre par fibre.
% slanted
\margin{vérifier cette affirmation !}
ou encore, en identifiant $q\mid D$ à un hom
\[
q^{\unsure{\prime}}:D^{\otimes2}\longrightarrow L
\]
que cet hom est un \uline{iso}
\[
q:D^{\otimes2}\xrightarrow{\sim}L\tag*{(6)}
\]
\deleted{La condition} En combinant avec l'iso. (5), on trouve un iso.
\[
L^{\otimes2}\simeq L
\]
i.e. un iso
\[
D\simeq L\simeq\mathcal O_S\tag*{(7)}
\]
Ce qu'on peut aussi exprimer ainsi :

\uline{Proposition} \inserted{Pour que $H=\mathbb P(V)$ soit non tangent à $Q\subset P$, il f et s} qu'\unsure{il|Il} existe une \deleted{unique} section $a$ de $T$ \deleted{$\unsure{V^\vee|V}$}, et une section $q$ du sous-fibré vectoriel $L^\vee$ de $\operatorname{Sym}^2(\check T)$ \unsure{définiss\textsuperscript{t}|définissant} $Q$, tels que

% margin note written sideways
\margin{Ceci équivaut aussi à : $q\mid\unsure{V^\perp|V}$ \uline{\unsure{régul.|régulière}}}
\[
\text{a)}\quad q(a)=1\tag*{(8)}
\]
\[
\text{b)}\quad\deleted{\unsure{E=\{|E=}}\pi(x)\unsure{=|\in}\deleted{\unsure{\Gamma(-,T)\mid|\Gamma(-,T)}}\varphi(a,x)\deleted{\unsure{=1\}}}\quad(x\in\Gamma(-,T)\unsure{)}
\tag*{(9)}
\]
\inserted{où $\varphi$ est la forme bil. associée à $q$.}

Alors $a,q$ sont déterminés de façon unique, et ce sont respectivement des bases de $D$ et de $L^\vee$. Enfin, on a
\[
\pi(a)=2\tag*{(10)}
\]
Il reste à prouver cette dernière relation. \deleted{Or \unsure{on a} plus gén, pour $x\in\Gamma(\unsure{-,|,}T)$}
% archive sheet number: 119
\scanpage{121}
\[
\deleted{\pi(x)=\varphi(a,x),}\tag*{\deleted{(11)}}
\]
Ce qu'on obtient en faisant $x=(a)$ dans (9), et notant que $\varphi(a,a)=2q(a)=2$

\inserted{\uline{Scholie} La cat des $(P,Q,H)$ comme dessus équivaut donc à celle des $(T,q,a)$, $T$ fibré vectoriel de \unsure{rg|rang} $2\nu$, $q$ forme quadratique \unsure{lisse} sur $T$, $a$ section de $T$, tels que $q(a)=1$}

\uline{NB} a) En car. résiduelles $\ne2$, on voit donc que $T\simeq V+D\simeq V+\mathcal O_S$ canoniquement, \inserted{(i.e. $E$ est muni d'une section origine canonique)} la décomposition étant ortho et $Q$ étant défini canoniquement par une forme
\[
q(x,t)=q_0(\unsure{x})+t^2\tag*{(11)}
\]
où $q_0$ est une forme quadratique sur $V$ ; $Q\cap E$ est défini par l'équation (en $x\in\Gamma(,E)$)
\[
q_0(x)+1=0\qquad\text{i.e.}\quad q_0(x)=-1\tag*{(12)}
\]
Dans les car. rés $\ne2$, la cat. des triples $(P,Q,H)$, ou des couples $(E,Q\cap E)$ ($E$ étant muni de sa structure de fibré affine \ldots) est \inserted{donc} équivalente à celle des couples $(V,q_0)$, où $V$ est un fibré vectoriel de rang $2\nu-1$.

En car. 2, par contre, on doit avoir automatiquement $D\subset V$, et il n'y a pas\deleted{\unsure{,}} de réduction au cas \unsure{d'une|$+$ une} forme de degré $2\nu-1$.

b) Si $T$ est de rang impair $2\nu+1$, donc $V$ de rang pair $2\nu$, les conditions que $Q$ soit
% archive sheet number: 120
\scanpage{122}
lisse s'expriment par la condition que
\[
\delta'(q)\in\Gamma(S,(\underline{\det}\,T)^{\otimes-2})\qquad\text{(discr. \uline{divisé})}\tag*{(13)}
\]
soit \uline{inv.}, et par la condition que
\[
D\xrightarrow{\pi\mid D}\mathcal O_S
\]
soit un iso, i.e.
\[
T\simeq V\oplus D\tag*{(14)}
\]
Ceci nous donne encore une section $a$ canonique de $D$, définie par la condition que
\[
\pi(a)=1\qquad\text{(non pas 2 !)}\tag*{(15)}
\]
et on peut alors \unsure{canoniser|normaliser} $q$ par la condition que
\[
q(a)=1\tag*{(16)}
\]
de sorte que l'on a encore
\[
q(x,t)=q_0(x)+t^2\tag*{(17)}
\]
et $Q$ se trouve défini par l'équation
\[
q_0(x)+1=0\qquad\text{i.e.}\quad q_0(x)=-1\qquad(x\in\Gamma(-,E))\tag*{(18)}
\]
où $E$ est identifié à $V$ grâce à (14).

Donc ici l'équivalence de cat. entre quadriques affines régulières \inserted{de rg $n$} et formes quadratiques \deleted{lisses} de rang $n-1$ marche,
% archive sheet number: 121
\scanpage{123}
sans restriction à la car. $\ne2$.

\uline{2} \quad On va définir une réflexion canonique dans $P$, ayant $H$ comme hyperplan de pts fixes, et invariant $Q$. Une \inserted{\unsure{pareille}} réflexion \deleted{\unsure{est associée} à un autom} ayant $H$ comme hyperplan fixe \deleted{s'écrit sous la forme} \inserted{induit $E$, donc est \unsure{associé|associée} à un autom. linéaire} $\unsure{\sigma|u}$ de la structure d'extension \deleted{$E$, \unsure{à la}} $T$ de $\mathcal O_S$ par $V$, \deleted{et} induisant l'identité sur $\mathcal O_S$, et une homothétie \inserted{$\lambda\,\mathrm{id}_V$} sur $V$, ou encore (quitte à diviser par $\lambda$\deleted{\unsure{\varphi}}) par un autom de la structure d'ext. induisant l'identité sur $V$, donc de la forme
\[
\sigma=\mathrm{id}+v
\]
où $v$ est nul sur $V$, donc de la forme
\[
v(x)=\pi(x).c\qquad\text{i.e.}\quad v=\pi\otimes c
\]
où $\pi\in\Gamma(S,\check T)$ est l'hom $T\longrightarrow\mathcal O_S$ définissant la structure d'extension, $c$ une section quelc. de $T$. Écrivons que \deleted{\unsure{Q}} $u$ invarie $q$ mod homothétie, on trouve
\[
\begin{aligned}
q\circ\sigma(x)&=q(x+vx)=q(x+\pi(x).c)\\
&=q(x)+\pi(x)^2q(c)+\deleted{\unsure{q}}\pi(x)\varphi(x,c)
\end{aligned}
\]
On veut trouver \deleted{\unsure{x}}$\lambda\in\Gamma(S,G_m)$ tel que
\[
\deleted{qu(x)=\lambda q(x)}\quad q\sigma(x)=\lambda q(x)\qquad\forall x\in\Gamma(-,T)
\]
% archive sheet number: 122
\scanpage{124}
i.e.
\[
(\lambda-1)q=\pi(q(c)\pi+\psi(c))\qquad
\left(\begin{aligned}\psi(c)&\in\Gamma(S,\check T)\\\psi(c)(x)&=\varphi(c,x)\end{aligned}\right)
\]
Comme $q$ est \deleted{non dégénérée}, lisse en tte fibre, ceci montre, si $\operatorname{rg}T\geq3$, i.e. $\operatorname{rg}E\geq2$, que $\lambda-1$ est nul en tte fibre (i.e que $\lambda-1$ doit être nilpotent), \deleted{\unsure{on}} et \deleted{que} comme $\pi$ est $\ne0$ en tte fibre, qu'on doit avoir $\psi(c)+q(c)\pi=0$ en tte fibre, \ldots\ Regardons les solutions correspondant à $\lambda=1$, i.e.
\[
\unsure{(|C}\qquad q\circ\sigma=q
\]
qui équivaut à
\[
q(c)\pi+\psi(c)=0\tag*{(19)}
\]
\deleted{On peut} On peut prendre p.ex. $c=-a$ (dans le cas \unsure{$T$|$E$} de rg pair, qui nous intéresse surtout)
\[
\sigma(x)=x\unsure{-|+}\pi(x)a\qquad(\operatorname{rg}\unsure{T|E}\ \text{pair})\tag*{(20)}
\]
% margin note written sideways
\margin{\uline{Prop} Il existe une unique réflexion de $P$ admettant $H$ comme hyperplan des pts fixes et \unsure{invariant|laissant} $Q$. Sa restriction à $\unsure{T|E}$ est donnée par (20) $\sigma x=x-\pi(x)a$ i.e. $\sigma=\mathrm{id}-\pi\otimes a$}
qui est de plus une réflexion i.e.
\[
\sigma^2=\mathrm{Id}_T\tag*{(21)}
\]
(car \deleted{\unsure{$\pi=$|$\sigma^2=$}} $\pi(a)=2\ldots$) En fait, c'est la seule solution \inserted{partt $\ne\mathrm{id}$}, \deleted{\unsure{on|dans}} car de (19) on tire $\psi(c)=\lambda\pi$ ($\lambda=-q(c)$) \deleted{\unsure{d'où}} i.e. $\psi(c)=\lambda\psi(a)=\psi(\lambda a)$ i.e. (comme $\psi$ est iso) $c=\underline{\lambda a}$, et il faut $\lambda=-q(c)$ i.e. $\lambda=-\lambda^2$ i.e. $\lambda(\lambda+1)=0$, ce qui (si on veut $\lambda$ partout non nul) donne $\lambda=-1$\ldots.

On dit que $\sigma$ est la \uline{symétrie canonique} de la quadrique affine régulière.
% archive sheet number: 123
\scanpage{125}
% letterhead
% upside down; the edge of the sheet cuts through the printed and handwritten lines
DES SCIENCES\\
\unsure{U} LANGUEDOC\\
\unsure{I}QUES

\margin{$\mathbb Z/2\cdot\unsure{\mathrm{ClP}(1)|\mathrm{CP}(1)|\mathbb C P(1)}\times\unsure{\mathrm{CP}(1|\mathrm{CO}(1|\mathbb C P(1)}$}
% margin note written sideways (upside down)

\[
\left\{\begin{aligned}
T&=M=\left\{\begin{pmatrix}a&b\\c&d\end{pmatrix}\right\}\\
q&=\deleted{\unsure{\mathrm{tr}}}\det=ad-bc
\end{aligned}\right.
\qquad\unsure{(b+b',(c+c')|(b+b',c+c')}
\]
\inserted{$\underline a=\underline1=\left(\begin{smallmatrix}1&0\\0&1\end{smallmatrix}\right)$}
\[
\deleted{\unsure{q(x)|\varphi(x)}}\ \varphi(a,b,c,d;a',b',c',d')
=b\underline{c'}+c\underline{b'}-(ad'+da')
\]
\[
\begin{aligned}
\pi(a,b,c,d)&=\varphi(a,b,c,d\unsure{;|,}1,0,0,1)=\unsure{+|-}(a+d)\\
&=\underline{-}\operatorname{Tr}\deleted{\unsure{\bigl\{}}\begin{pmatrix}a&b\\c&d\end{pmatrix}
\end{aligned}
\]
\inserted{attention au signe !}
\[
\deleted{V}\qquad V=\unsure{\gamma|\mathfrak g}M=\left\{\begin{pmatrix}a&b\\c&-a\end{pmatrix}\right\}
\]
% archive sheet number: 124
\scanpage{126}
Si $T$ est de rang impair, donc $T\simeq V\oplus\mathcal O_S\ldots$, on définira \deleted{\unsure{une}} $\sigma$ par
\[
\sigma(x)=x-2\pi(x)a\tag*{(22)}
\]
ou encore, sur $E$, par la symétrie p.r. à l'origine $\unsure{a|0}$. On aura encore (21).

3. \uline{Schéma des autom d'une quadrique affine} \inserted{birégulière}.

Soit
\[
\begin{aligned}
G&=\underline{\operatorname{Aut}}(P,H,Q)=\underline{\operatorname{Aut}}\,\mathrm{aff}(E,Q\cap E)\\
&\deleted{\simeq\operatorname{Aut}(T,\unsure{q},\pi\unsure{)},a,q)}
\end{aligned}
\]
\[
\deleted{(\pi\in\Gamma(S,\check T),a\in\Gamma(S,T),q(a)=1,\varphi(a,x)=}
\]
Dans le cas où $T$ de rang impair i.e. $E$ de rang pair, on a
\[
G\simeq\operatorname{Aut}(V,q_0)\qquad V\ \text{de rang pair}
\]
et désignant par $G^+$ le sous-schéma en groupes qui opère trivialement sur $\det V$, on trouve
\[
\deleted{\unsure{G\simeq G^+\times\mu_2|G\simeq G^+\times\mathbb Z/2\mathbb Z}}
\]
\diagram{\[
\begin{tikzcd}[ampersand replacement=\&,column sep=large]
\unsure{0|\cdot|1} \arrow[r,phantom,"{\unsure{-|\longrightarrow}}" description] \& G^+ \arrow[r] \& G \arrow[r] \& \mathbb Z/2\mathbb Z \arrow[r,phantom,"{\unsure{-|\longrightarrow}}" description] \& 0
\end{tikzcd}
\]}
Si $\unsure{T|E}$ de rang pair, \inserted{$G\simeq\operatorname{Aut}(\unsure{T|E},q,a)$} on va montrer qu'on a un iso can.
\[
G\simeq G^0\times(\mathbb Z/2\mathbb Z)_S
\]
où $G^0\xrightarrow{\sim}\underline{\operatorname{Aut}}^+(V,q_0)$, $(\mathbb Z/2\mathbb Z)_S$ engendré par $\sigma$.
% archive sheet number: 125
\scanpage{127}
(31 p)

Formes quadratiques

Fibrés en drapeaux quad, groupes orthogonaux
% archive sheet number: 126
\scanpage{128}
Formes quadratiques, quadriques et groupes orthogonaux

1. Discriminant et discriminant divisé -- formes quadratiques régulières (ou lisses)

Si $q$ est une forme quadratique sur $E$\inserted{ de rang $n$}, on a discr. div.
\[
\delta'(q)\in\Gamma(\det\check E)^{\otimes2}=\Gamma(\det E)^{\otimes(-2)}\tag*{(1.1)}
\]
\margin{(qui est le disc. ordinaire si $n$ pair)}
\margin{\uline{NB} $\delta'(\lambda q)=\lambda^{\unsure{n|h}}\delta'(q)$}
Conditions équivalentes sur $q$.

a) $\delta'(q)$ inversible

b) $V(q)\subset\check{\mathbb P}(E)$ est lisse, de codim 1 en chaque fibre (ou encore : est un diviseur relatif \uline{lisse})

c) $\forall s\in S$, $q\otimes k(s)$\deleted{\unsure{$\in$}} appartient à l'unique orbite ouverte de $\underline{\mathrm{Gl}}(E(s))$ opérant sur le \deleted{\unsure{lieu}}\inserted{schéma} $\underline{\mathrm{Quad}}(E(s))$ (qui est aussi l'unique orbite de dimension maximale, laquelle dimension est $\frac{n(n+1)}{2}$)

d) $\forall s\in S$, $\underline{\operatorname{Aut}}(q\otimes k(s))$ est de dimension minimale (savoir $\frac{n(n-1)}{2}$).

\unsure{c'|e'|d'}) \inserted{$\forall s\in S$,} $\exists$ ouvert \deleted{$\unsure{\varnothing\ne}$}\inserted{non vide} $U\subset\underline{\mathrm{Quad}}(E(s))$ tel que $q\otimes k(s)$ équiv\textsuperscript{\unsure{lt|te|alente}} (sur ext. conv. de $k(s)$) aux formes $\in U$ [NB on peut prendre $U$ défini par $\delta'$ considéré comme fonction pol. sur $\underline{\mathrm{Quad}}(E)\ldots$]

e) $q$ est \unsure{isom.} loc! fppf [\deleted{ou} aussi : loc : (ét.) si $2$ inv. sur $S$ \uline{ou} $n$ pair] à la forme \unsure{standart|standard} $q_n$

\uline{NB} Si $q$ est à valeurs dans Mod. inv. $L$, alors forme bil. \inserted{\unsure{ass.}} $\varphi$ est à val. dans $L$, d'où $E\unsure{-|\longrightarrow}\check E\otimes L$, d'où $\det E\unsure{-|\longrightarrow}\det(\check E\otimes L)\simeq(\det\check E)\otimes L^{\otimes n}$, donc
\[
\delta'(q)\in\Gamma{\unsure{\det(E)|(\det E)}}^{\otimes(-2)}\otimes L^{\otimes n}=\operatorname{Hom}((\det E)^{\otimes2},L^{\otimes n})\tag*{(1.2)}
\]
% archive sheet number: 127
